\documentclass[11pt]{article}

\usepackage[margin=1in]{geometry}
\usepackage{amsmath,amssymb,amsthm,mathtools}
\usepackage{bm}
\usepackage{booktabs}
\usepackage{enumitem}
\usepackage{microtype}
\usepackage{xcolor}
\usepackage{hyperref}
\usepackage[nameinlink,capitalise]{cleveref}

\hypersetup{
  colorlinks=true,
  linkcolor=blue!55!black,
  citecolor=blue!55!black,
  urlcolor=blue!55!black,
  pdftitle={Causal Capacity Quantum Gravity},
  pdfauthor={Payton Ison}
}

\newtheorem{definition}{Definition}[section]
\newtheorem{axiom}[definition]{Axiom}
\newtheorem{conjecture}[definition]{Conjecture}
\newtheorem{proposition}[definition]{Proposition}
\newtheorem{theorem}[definition]{Theorem}
\newtheorem{lemma}[definition]{Lemma}
\newtheorem{corollary}[definition]{Corollary}
\theoremstyle{remark}
\newtheorem{remark}[definition]{Remark}
\newtheorem{obligation}[definition]{Proof obligation}

\newcommand{\cH}{\mathcal H}
\newcommand{\cA}{\mathcal A}
\newcommand{\cC}{\mathcal C}
\newcommand{\cK}{\mathcal K}
\newcommand{\cW}{\mathcal W}
\newcommand{\cG}{\mathcal G}
\newcommand{\cD}{\mathcal D}
\newcommand{\Tr}{\operatorname{Tr}}
\newcommand{\id}{\mathbf 1}
\newcommand{\dd}{\mathrm d}
\newcommand{\order}{\mathcal O}
\newcommand{\mpl}{M_{\mathrm P}}
\newcommand{\ellp}{\ell_{\mathrm P}}
\newcommand{\gsm}{G_{\mathrm{SM}}}
\newcommand{\doi}[1]{\href{https://doi.org/#1}{doi:#1}}

\title{\textbf{Causal Capacity Quantum Gravity}\\
\large A Background-Free Information-Theoretic Bootstrap for
Four-Dimensional Semiclassical Gravity}
\author{Payton Ison}
\date{\today}

\begin{document}
\maketitle

\begin{abstract}
We formulate a finite-cutoff, background-free quantum-gravity research
program in which causal composition and quantum information capacity are
microscopic, while Lorentzian geometry is emergent.  A cutoff theory is a
finite direct sum of relational graph states.  Its elementary transitions
are local graph moves generated by a self-adjoint Hamiltonian, making
finite-cutoff evolution exactly unitary.  Every causal cut carries a
central capacity charge \(s=\log q\); in a semiclassical phase this charge
is identified with area through \(A=4G\hbar s\).  Discrete bivector data and
simplicity, closure, and shape-matching constraints supply the geometric
information that areas alone cannot contain.

We state a distributional refinement prescription and isolate five
nontrivial continuum requirements: a non-topological fixed point, an
anomaly-free hypersurface-deformation algebra, a single Lorentzian metric
with exactly two gapless graviton helicities, a chiral Standard Model
sector, and a finite unitary continuation through high-curvature
black-hole states.  Conditional on these requirements and a
code-subspace modular identity, we prove that the two-derivative infrared
theory is semiclassical Einstein gravity universally coupled to quantum
matter.  The construction is therefore not presented as a completed
theory of quantum gravity.  It is a precise bootstrap: finite kinematics
and dynamics are specified, the infrared closure is established
conditionally, and the remaining existence claims are exposed as
falsifiable mathematical proof obligations.
\end{abstract}

\tableofcontents

\section{Introduction}

The two successful dynamical frameworks of modern physics begin from
apparently incompatible primitives.  Quantum theory assigns states,
operator algebras, and amplitudes to experiments.  General relativity
assigns causal and metrical relations to events, while making the metric
itself dynamical.  Black-hole thermodynamics suggests that neither set of
primitives is complete in isolation: a surface of area \(A\) carries an
entropy \(A/(4G\hbar)\), and semiclassical gravitational field equations
can be related to local horizon thermodynamics and entanglement
equilibrium \cite{Bekenstein1973,Hawking1975,Jacobson1995,Jacobson2015}.

Entropy by itself, however, cannot be fundamental.  Distinct pure states
all have vanishing von Neumann entropy, and states with the same entropy
can have different phases, correlations, and operational content.  The
more faithful primitive is \emph{quantum information capacity}: the
dimension of an admissible state space, together with its operator
algebra, causal composition law, and dynamics.  Entropy is a
state-dependent use of that capacity.

This paper develops that distinction into a candidate architecture called
\emph{Causal Capacity Quantum Gravity} (CCQG).  It combines four ideas:
\begin{enumerate}[label=(\roman*)]
  \item locally finite quantum causal histories
  \cite{Markopoulou2000,Hawkins2003};
  \item dynamical combinatorial geometry and covariant amplitudes
  \cite{Regge1961,RovelliVidotto2015};
  \item generalized entropy and modular reconstruction
  \cite{Jacobson2015,JLMS2016,ADH2015};
  \item a non-topological, distributional continuum limit
  \cite{Bruno2026}.
\end{enumerate}

The main result is deliberately conditional.  We give a finite
microscopic ansatz and prove that, if its continuum limit satisfies a
short list of sharply stated conditions, its infrared dynamics closes on
Einstein gravity coupled universally to quantum matter.  We do
\emph{not} prove that the required continuum fixed point exists, that the
Standard Model is uniquely selected, or that the black-hole continuation
survives removal of the cutoff.  Those are the central proof obligations
of the program.

\subsection{Claim ledger}

For clarity, the logical status of the principal statements is:
\begin{center}
\begin{tabular}{p{0.31\linewidth}p{0.58\linewidth}}
\toprule
\textbf{Status} & \textbf{Statement}\\
\midrule
Definition &
Finite relational state space, capacity labels, local-move Hamiltonian,
and regulated transition amplitude.\\
Exact finite-cutoff result &
Evolution generated by the regulated self-adjoint Hamiltonian is unitary.\\
Semiclassical matching condition &
Capacity becomes area and the large-label phase becomes the Regge action.\\
Conditional theorem &
An anomaly-free single-metric continuum phase has Einstein gravity as its
two-derivative infrared theory.\\
Open existence conjecture &
The regulated models possess the required four-dimensional,
non-topological continuum limit with chiral matter and regular black-hole
evolution.\\
\bottomrule
\end{tabular}
\end{center}

\section{Microscopic kinematics}

\subsection{Relational causal complexes}

Let the regulator be \(\Lambda=(N,Q)\), where \(N\) bounds the number of
events and \(Q\) bounds local information capacity.  Let
\(\mathfrak G_\Lambda\) be a finite set of oriented, locally finite,
acyclic two-complexes \(K\), considered modulo combinatorial
automorphisms.  Orientation determines a partial order on events, but no
metric, embedding manifold, coordinate system, or preferred lattice
spacing is supplied.

A spatial cut \(\Gamma\) is an antichain intersecting every causal path in
the portion of \(K\) under consideration.  Each face \(f\) crossing
\(\Gamma\) has a sector label
\begin{equation}
  \zeta_f=(q_f,\beta_f,r_f), \qquad
  1\le q_f\le Q.
\end{equation}
Here \(q_f\) is an information-capacity label, \(\beta_f\) is discrete
geometric data, and \(r_f\) labels an internal gauge representation.
The allowed \(\beta_f\) belong to a finite regulator set
\(\mathcal B_Q\).

\begin{definition}[Capacity charge]
The capacity charge of a face is
\begin{equation}
  s_f=\log q_f.
\end{equation}
For a cut \(\gamma\), the additive capacity operator is the central
operator
\begin{equation}
  \widehat s_\gamma
  =\bigoplus_{\bm\zeta}
   \left(\sum_{f\subset\gamma}\log q_f\right)
   \id_{\bm\zeta}.
  \label{eq:capacity}
\end{equation}
\end{definition}

The use of \(\log q\) is forced by additivity under tensor composition:
\(\log(q_1q_2)=\log q_1+\log q_2\).  It should not be confused with the
entropy of a particular state.  In a fixed sector,
\begin{equation}
  S(\rho_\gamma)\le \log\dim\cH_{\gamma,\bm\zeta},
\end{equation}
with equality only for a maximally mixed state.  If superselection sectors
are present, the entropy also contains a classical sector term
\begin{equation}
  S(\rho_\gamma)
  =H(\{p_{\bm\zeta}\})
   +\sum_{\bm\zeta}p_{\bm\zeta}S(\rho_{\gamma,\bm\zeta}).
\end{equation}
Thus area will be associated with a central capacity charge, not with an
arbitrary von Neumann entropy.

\subsection{Boundary Hilbert spaces and nonfactorization}

For a cut \(\Gamma\), define
\begin{equation}
  \cH_{\Gamma,\Lambda}
  =
  \bigoplus_{\bm\zeta}
  \left[
    \cH^{\mathrm{geom}}_{\Gamma,\bm\zeta}
    \otimes
    \cH^{\mathrm{matter}}_{\Gamma,\bm\zeta}
  \right],
  \label{eq:cutHilbert}
\end{equation}
and the regulated kinematical space
\begin{equation}
  \cH_\Lambda
  =
  \bigoplus_{\Gamma\in\mathfrak G_\Lambda}
  \cH_{\Gamma,\Lambda}\big/\mathrm{Aut}(\Gamma).
  \label{eq:totalHilbert}
\end{equation}
Because \(\mathfrak G_\Lambda\), \(\mathcal B_Q\), and all local state
spaces are finite, \(\cH_\Lambda\) is finite dimensional.

Equation \eqref{eq:cutHilbert} is a direct sum over shared central sectors,
not a universal factorization into independent ``inside'' and ``screen''
Hilbert spaces.  This distinction is required by gauge constraints and is
also the natural structure of operator-algebra quantum error correction
\cite{ADH2015,JLMS2016}.

\subsection{Geometric constraints}

Capacity data alone do not determine a geometry: different tetrahedra can
have identical face areas.  The label \(\beta_f\) therefore carries a
regulated version of an oriented bivector \(B_f^{IJ}\).  Semiclassical
geometric sectors are those satisfying, up to regulator errors,
\begin{align}
 n_{eI}B_f^{IJ}&=0
 &&\text{(simplicity)}, \label{eq:simplicity}\\
 \sum_{f\subset e}\epsilon_{ef}B_f^{IJ}&=0
 &&\text{(closure)}, \label{eq:closure}
\end{align}
together with shape matching across adjacent cells and a nondegeneracy
condition.  These conditions reconstruct a piecewise-flat Lorentzian
geometry in the large-label regime.

\begin{axiom}[Capacity--area matching]
There is a universal microscopic area-per-nat constant \(a_\star\) such
that, on semiclassical geometric states,
\begin{equation}
  \widehat A_\gamma=a_\star\widehat s_\gamma
  +\order(Q^{-1}).
  \label{eq:areaCapacity}
\end{equation}
Only after identifying the infrared gravitational coupling do we set
\begin{equation}
  a_\star=4G\hbar.
  \label{eq:aStar}
\end{equation}
\end{axiom}

Keeping \(a_\star\) independent in the microscopic definition avoids
assuming Newton's constant before it has an infrared operational meaning.

\section{Finite background-free dynamics}

\subsection{Local moves}

Let \(m:\Gamma\rightarrow\Gamma'\) be an allowed local causal move,
including graph refinement, coarse graining, reconnection, and local
matter evolution.  Denote its partial isometry by
\(\widehat M_{m,\Gamma}\).  Define the regulated move Hamiltonian
\begin{equation}
  \widehat H_\Lambda
  =
  \widehat H_{\mathrm{diag}}
  +\sum_{\Gamma,m}
   \left(
     \lambda_{m,\Gamma}\widehat M_{m,\Gamma}
     +\overline{\lambda}_{m,\Gamma}
      \widehat M_{m,\Gamma}^{\dagger}
   \right),
  \label{eq:moveHamiltonian}
\end{equation}
where \(\widehat H_{\mathrm{diag}}\) is self-adjoint and invariant under
graph automorphisms.  The coefficients may depend only on relational
labels in the causal neighborhood of the move.  No fixed geometric
distance occurs in \eqref{eq:moveHamiltonian}.

A clock degree of freedom \(T\) may be introduced relationally, with
constraint
\begin{equation}
  \widehat C_\Lambda
  =
  \widehat p_T+\widehat H_\Lambda.
\end{equation}
Conditional evolution with respect to \(T\) is
\begin{equation}
  U_\Lambda(\tau)
  =\exp\left(-\frac{i}{\hbar}\tau\widehat H_\Lambda\right),
\end{equation}
and the finite transition amplitude is
\begin{equation}
  W_\Lambda(y_+,y_-;\tau)
  =
  \langle y_+|U_\Lambda(\tau)|y_-\rangle.
  \label{eq:finiteAmplitude}
\end{equation}
Equivalently, group averaging of \(\widehat C_\Lambda\) defines a
regulated rigging map.

\begin{proposition}[Finite-cutoff unitarity]
For every finite regulator \(\Lambda\), the amplitude
\eqref{eq:finiteAmplitude} defines exact unitary relational evolution.
\end{proposition}

\begin{proof}
The space \(\cH_\Lambda\) is finite dimensional and
\(\widehat H_\Lambda=\widehat H_\Lambda^\dagger\) by
\eqref{eq:moveHamiltonian}.  The spectral theorem therefore gives
\(U_\Lambda(\tau)^\dagger U_\Lambda(\tau)=\id\) for every real \(\tau\).
\end{proof}

\begin{remark}
A complex sum of completely positive maps is not generally a quantum
channel.  CCQG therefore sums amplitudes through a self-adjoint generator;
completely positive maps arise only after restricting or tracing a global
unitary evolution.
\end{remark}

\subsection{Semiclassical phase condition}

The microscopic dynamics should not insert a background Einstein metric.
It must nevertheless possess an asymptotic geometric phase.  For
large-capacity coherent boundary states satisfying
\eqref{eq:simplicity}--\eqref{eq:closure}, require
\begin{equation}
\begin{split}
  W_\Lambda
  \sim
  \sum_{\text{saddles }g}
  \mathcal N_g
  \exp\Bigg\{i\Bigg[
    \frac{1}{2\pi}\sum_f s_f\epsilon_f
    -\frac{\lambda_{\mathrm{cc}}}{8\pi G\hbar}
       \sum_\sigma V_\sigma\\
    +\frac{1}{\hbar}I_{\mathrm{matter}}[g,\Phi]
    +\order\!\left(\frac{1}{s},\ell_\star^2R^2\right)
  \Bigg]\Bigg\}.
  \label{eq:ReggeAsymptotic}
\end{split}
\end{equation}
Here \(\epsilon_f\) is the Regge deficit angle, \(V_\sigma\) is a
four-volume, and \(\lambda_{\mathrm{cc}}\) is the cosmological constant.
Using \eqref{eq:aStar}, the first term is precisely
\(I_{\mathrm{Regge}}/\hbar\):
\begin{equation}
  \frac{1}{8\pi G\hbar}\sum_f A_f\epsilon_f
  =\frac{1}{2\pi}\sum_f s_f\epsilon_f.
\end{equation}

This matching condition is not itself a derivation of gravity.  It is the
microscopic universality condition whose emergence under renormalization
must be proved.  The canonical conjugate of a boundary area charge is a
boundary boost angle; a useful semiclassical normalization is
\begin{equation}
  [\widehat\eta_f,\widehat s_{f'}]
  =2\pi i\,\delta_{ff'}+\order(Q^{-1}).
  \label{eq:cornerCanonical}
\end{equation}
The bulk deficit angle should not be confused with this boundary corner
variable.

\section{Refinement and the continuum}

Let \(\Lambda\prec\Lambda'\) denote removal of a regulator restriction,
and let
\(\iota_{\Lambda\Lambda'}:\cH_\Lambda\rightarrow\cH_{\Lambda'}\)
be an isometric refinement map.  Exact invariance under all local
refinements would risk erasing propagating degrees of freedom and
producing a topological theory.  Recent model-independent analysis of
spin-foam limits makes this obstruction precise under strong convergence
assumptions \cite{Bruno2026}.  We therefore require neither exact Pachner
invariance nor norm convergence of every regulated amplitude.

\begin{definition}[Distributional cylindrical limit]
A continuum amplitude exists if there is a dense test space
\(\mathcal D\), a physical distribution space \(\mathcal D^\ast\), and a
rigging map \(\eta:\mathcal D\rightarrow\mathcal D^\ast\) such that, for
all test states \(\psi,\phi\),
\begin{equation}
  \lim_{\Lambda\rightarrow\infty}
  \langle
    \iota_\Lambda\psi,
    U_\Lambda\,\iota_\Lambda\phi
  \rangle
  =
  \eta(\psi)[\phi],
  \label{eq:weakLimit}
\end{equation}
after renormalization of the move couplings.  The limit is required only
in the weak topology selected by physical observables.
\end{definition}

\begin{conjecture}[Non-topological refinement fixed point]
There exists a non-Gaussian renormalization trajectory
\(\lambda_{m,\Gamma}(\Lambda)\) for which
\eqref{eq:weakLimit} exists, carries local propagating excitations, and is
independent of regulator details in its infrared observables.
\label{conj:fixedPoint}
\end{conjecture}

The continuum theory must recover type-III local operator algebras even
though each regulated algebra is finite dimensional.  Finite
factorization is therefore a regulator property, not a claim about exact
continuum subregions.

\section{Modular matching}

Let \(V:\cH_{\mathrm{code}}\rightarrow\cH_R\) encode a semiclassical bulk
code subspace into the algebra of an accessible region \(R\).  For a
reference state \(\sigma\), write
\(K^\sigma=-\log\sigma\) for its modular Hamiltonian.

\begin{axiom}[Code-subspace modular identity]
For every sufficiently large semiclassical causal diamond,
\begin{equation}
\begin{split}
 P_{\mathrm{code}}V^\dagger K_R^\sigma V P_{\mathrm{code}}
 =
 P_{\mathrm{code}}\left(
   \frac{\widehat A_\gamma}{4G\hbar}
   +K_r^\sigma+K_{\mathrm{ct}}
 \right)P_{\mathrm{code}}
 +c_\sigma P_{\mathrm{code}}
 +\order(\ell_\star/L).
 \label{eq:modularIdentity}
\end{split}
\end{equation}
\end{axiom}

This is a reference-state, code-subspace statement modeled on the
bulk--boundary relative-entropy relation in holography
\cite{JLMS2016}.  It is not imposed as a nonlinear operator identity on
the full microscopic Hilbert space.

Equation \eqref{eq:modularIdentity} identifies the generalized entropy
\begin{equation}
  S_{\mathrm{gen}}(\gamma)
  =
  \frac{\langle A_\gamma\rangle}{4G\hbar}
  +S_{\mathrm{bulk}}+S_{\mathrm{ct}}.
  \label{eq:Sgen}
\end{equation}
The quantum first law,
\begin{equation}
  \delta S_{\mathrm{bulk}}
  =\delta\langle K^\sigma_{\mathrm{bulk}}\rangle,
  \label{eq:firstLaw}
\end{equation}
then converts state variations into geometric response.

\section{Conditional infrared closure}

We now state the part of the program that can be closed without assuming a
specific microscopic fixed point.

\begin{theorem}[Conditional infrared Einstein limit]
Suppose the continuum limit of CCQG has:
\begin{enumerate}[label=(\roman*)]
  \item one \(3+1\)-dimensional Lorentzian metric and a common causal cone;
  \item local Lorentz invariance, unitarity, cluster decomposition, and a
        derivative expansion;
  \item no gapless geometric excitation other than a rank-two mode;
  \item an anomaly-free representation of the hypersurface-deformation
        algebra; and
  \item the modular identity \eqref{eq:modularIdentity} and stationary
        generalized entropy for all sufficiently small vacuum diamonds.
\end{enumerate}
Then the two-derivative infrared field equation is
\begin{equation}
  G_{\mu\nu}+\lambda_{\mathrm{cc}}g_{\mu\nu}
  =
  8\pi G\,\langle T_{\mu\nu}\rangle,
  \label{eq:semiclassicalEinstein}
\end{equation}
up to higher-derivative and nonlocal effective-field-theory corrections.
\end{theorem}

\begin{proof}
Assumption (iv) gives, in continuum notation,
\begin{align}
 \{D[\bm N],D[\bm M]\}
   &=D[[\bm N,\bm M]],\\
 \{D[\bm N],H[M]\}
   &=H[\mathcal L_{\bm N}M],\\
 \{H[N],H[M]\}
   &=D[h^{ij}(N\partial_jM-M\partial_jN)].
 \label{eq:HDA}
\end{align}
Under the locality and embeddability assumptions of the
Hojman--Kucha\v{r}--Teitelboim reconstruction, this algebra determines the
ADM form of metric dynamics up to coupling constants and a cosmological
term \cite{HKT1976}.  The six components of \(h_{ij}\) give twelve
phase-space variables.  Four first-class constraints remove eight
phase-space dimensions, leaving four phase-space, or two configuration,
degrees of freedom: the two graviton helicities.

Linearizing about a semiclassical background yields the gauge symmetry
\begin{equation}
  \delta h_{\mu\nu}
  =\nabla_\mu\xi_\nu+\nabla_\nu\xi_\mu.
\end{equation}
Lorentz invariance of amplitudes involving a soft massless spin-two
particle forces a universal coupling to the conserved stress tensor
\cite{Weinberg1964}.  Since the spin-two field itself carries energy,
consistency iterates this coupling into the nonlinear Einstein
self-interaction \cite{Deser1970}.  Consequently the most general
infrared effective action begins
\begin{equation}
\begin{split}
 I_{\mathrm{IR}}
 =
 \int\dd^4x\sqrt{-g}\Big[
   \frac{1}{16\pi G}(R-2\lambda_{\mathrm{cc}})
   +\mathcal L_{\mathrm{matter}}
   +c_1R^2+c_2R_{\mu\nu}R^{\mu\nu}
   +\cdots
 \Big].
 \label{eq:IREFT}
\end{split}
\end{equation}

Independently, apply \eqref{eq:firstLaw} and stationary
\(S_{\mathrm{gen}}\) to every small causal diamond.  Raychaudhuri
focusing relates the area variation to
\(R_{\mu\nu}k^\mu k^\nu\), while the local modular Hamiltonian relates
the matter variation to
\(\langle T_{\mu\nu}\rangle k^\mu k^\nu\).  Therefore
\begin{equation}
  \left(
    G_{\mu\nu}
    -8\pi G\langle T_{\mu\nu}\rangle
  \right)k^\mu k^\nu=0
  \label{eq:nullEquation}
\end{equation}
for every null vector \(k^\mu\) \cite{Jacobson1995,Jacobson2015}.
The tensor in parentheses must be proportional to \(g_{\mu\nu}\).
The Bianchi identity and stress-energy conservation make the
proportionality constant spacetime independent, yielding
\eqref{eq:semiclassicalEinstein}.  Higher-curvature terms in
\eqref{eq:IREFT} correspond to Wald-like corrections to generalized
entropy \cite{Wald1993,Bueno2017}.
\end{proof}

\begin{remark}[Weinberg--Witten]
An ordinary local quantum field theory on a fixed background cannot
generically produce a composite massless spin-two particle carrying a
Lorentz-covariant conserved stress tensor \cite{WeinbergWitten1980}.
CCQG evades the premise only if its substrate is genuinely
background-free and the gravitational stress tensor is gauge dependent,
as in general relativity.  This is a constraint on the construction, not
an automatic loophole.
\end{remark}

\section{The Standard Model sector}

Reproducing the Standard Model is distinct from deriving it uniquely.
The minimal present claim is an embedding.

\begin{axiom}[Internal algebra and representation content]
The matter labels admit the gauge group
\begin{equation}
  \gsm
  =
  \frac{SU(3)_c\times SU(2)_L\times U(1)_Y}{\mathbb Z_6},
\end{equation}
with three copies of the chiral representation
\begin{equation}
\begin{array}{c|ccccc}
 & Q_L & u_R & d_R & L_L & e_R\\
\hline
SU(3)_c\times SU(2)_L\times U(1)_Y
&(3,2)_{1/6}&(3,1)_{2/3}&(3,1)_{-1/3}
&(1,2)_{-1/2}&(1,1)_{-1}
\end{array}
\end{equation}
and a Higgs sector \((1,2)_{1/2}\).  A sterile right-handed neutrino may
be included without altering gauge anomalies.
\end{axiom}

Gauge holonomies live on relational links, curvature on faces, and
fermions on events or dual cells.  The local-move Hamiltonian contains
gauge-invariant Wilson terms, a regulated chiral Dirac operator, Yukawa
couplings, and the Higgs potential.  The chiral operator is required to
satisfy a Ginsparg--Wilson relation,
\begin{equation}
  \gamma_5D+D\gamma_5
  =aD\gamma_5D,
  \label{eq:GW}
\end{equation}
or its background-free analogue.  This is the known route around naive
fermion doubling while preserving the correct anomaly structure
\cite{GinspargWilson1982,Luscher1999}.

The usual local anomaly constraints must hold, including
\begin{equation}
 [SU(3)]^2U(1),\quad
 [SU(2)]^2U(1),\quad
 [U(1)]^3,\quad
 [\mathrm{gravity}]^2U(1),
\end{equation}
together with the absence of the global \(SU(2)\) anomaly.  The listed
Standard Model multiplets satisfy these constraints generation by
generation.

\begin{obligation}[Chiral continuum matter]
Construct the analogue of \eqref{eq:GW} on the fluctuating causal
complexes, prove locality relative to the emergent causal structure, and
show that the weak refinement limit contains precisely the stated chiral
spectrum with no mirror fermions.
\label{obl:chiral}
\end{obligation}

The gauge group, three generations, Yukawa matrices, and Higgs parameters
are inputs at this stage.  A stronger theory would derive them from the
finite internal algebra or from stability of the refinement fixed point.
Until that derivation exists, CCQG reproduces rather than explains the
Standard Model.

\section{Black holes and the no-geometry phase}

\subsection{Finite-cutoff continuation}

At finite \(\Lambda\), every relational observable is a bounded operator
on a finite-dimensional Hilbert space and \(U_\Lambda(\tau)\) exists for
all real \(\tau\).  A classical curvature singularity therefore cannot be
a failure of the microscopic evolution law.  Instead it marks failure of
the semiclassical geometric constraints
\eqref{eq:simplicity}--\eqref{eq:closure} or of the saddle expansion
\eqref{eq:ReggeAsymptotic}.  The state can enter a causal quantum phase in
which no nondegenerate metric is defined.

This observation establishes regulated continuation, not continuum
singularity resolution.  The latter requires the limiting relational
amplitudes to remain finite and to connect appropriate semiclassical
in- and out-states.

\subsection{Encoding and evaporation}

A horizon cut carries the capacity
\begin{equation}
  s_{\mathrm{BH}}
  =\frac{A_{\mathrm H}}{4G\hbar},
\end{equation}
but the microscopic state need not be maximally mixed.  Interior
operators are logical operators encoded nonlocally in horizon and
radiation algebras, as suggested by holographic quantum error correction
\cite{ADH2015,Pastawski2015}.

At relational time \(\tau\), global evaporation is an isometry
\begin{equation}
  V_\tau:\cH_{\mathrm{in}}
  \longrightarrow
  \cH_{\mathrm{BH}}(\tau)\otimes\cH_{\mathrm{rad}}(\tau).
  \label{eq:evapIsometry}
\end{equation}
Radiation alone evolves by the completely positive map
\begin{equation}
  \mathcal E_\tau(\rho)
  =
  \Tr_{\mathrm{BH}}\!\left(
    V_\tau\rho V_\tau^\dagger
  \right).
\end{equation}
If the final black-hole factor becomes one dimensional, the final
radiation map is itself an isometry and global information is recovered.
Quantum extremal surfaces extremize \eqref{eq:Sgen}
\cite{EngelhardtWall2015}; island calculations provide a semiclassical
mechanism for the associated Page transition
\cite{Penington2020,Almheiri2020}.

\begin{conjecture}[Unitary high-curvature continuation]
There is a continuum physical channel obtained from
\eqref{eq:evapIsometry} whose relational amplitudes remain finite through
the no-geometry phase, whose asymptotic in- and out-spaces are complete,
and whose final inaccessible factor is trivial.
\label{conj:blackHole}
\end{conjecture}

Finite-dimensional unitarity alone does not prove
\cref{conj:blackHole}; the order of the large-system, continuum, and
late-time limits matters.

\section{Central existence conjecture}

The entire program can be compressed into one statement.

\begin{conjecture}[CCQG existence conjecture]
There exists a sequence of finite relational models
\((\cH_\Lambda,\widehat H_\Lambda,\iota_{\Lambda\Lambda'})\) such that:
\begin{enumerate}[label=(\arabic*)]
  \item each regulated amplitude is finite, background-free, and exactly
        unitary;
  \item the weak refinement limit \eqref{eq:weakLimit} exists and is
        non-topological;
  \item the continuum constraint algebra is anomaly free;
  \item the infrared spectrum contains exactly one four-dimensional
        Lorentzian metric and two graviton helicities;
  \item the modular relation \eqref{eq:modularIdentity} holds in every
        semiclassical code subspace;
  \item the chiral matter limit is the Standard Model sector described
        above; and
  \item black-hole states satisfy
        \cref{conj:blackHole}.
\end{enumerate}
\label{conj:master}
\end{conjecture}

\Cref{conj:master} is stronger than a unification slogan: every clause
can fail independently in a calculation.  It is also weaker than a
completed theory, because it is an existence statement rather than an
explicit fixed-point solution.

\section{Proof program and falsification criteria}

The next mathematical work is naturally divided into five theorems.

\begin{obligation}[Finite regulator and measure]
Choose the allowed move set and couplings, prove essential
self-adjointness before the finite cutoff is taken, and show that gauge
projection does not introduce negative-norm states.
\end{obligation}

\begin{obligation}[Non-topological continuum]
Construct the refinement maps and renormalization flow, establish the
distributional limit \eqref{eq:weakLimit}, and exhibit a nonzero
two-point function for a propagating geometric excitation.
\end{obligation}

\begin{obligation}[Four-dimensional gravitational phase]
Show that the long-distance spectral and Hausdorff dimensions approach
four, derive a unique causal cone, prove the hypersurface-deformation
algebra without anomaly, and demonstrate that all unwanted geometric
modes are gapped.
\end{obligation}

\begin{obligation}[Modular emergence]
Derive \eqref{eq:modularIdentity} from the microscopic operator algebra,
including edge-mode and counterterm contributions, rather than imposing
it on semiclassical states.
\end{obligation}

\begin{obligation}[Matter and black-hole completion]
Prove \cref{obl:chiral,conj:blackHole}, including the order of limits
needed for a unitary asymptotic scattering map.
\end{obligation}

The framework is falsified, in its present form, by any of the following:
\begin{enumerate}
  \item every refinement fixed point is topological or lacks a positive
        physical inner product;
  \item the continuum constraint commutators contain an incurable anomaly;
  \item the infrared phase contains extra gapless scalar, vector, or
        spin-two modes;
  \item capacity--area matching is incompatible with the semiclassical
        graviton normalization;
  \item chiral gauge matter cannot be defined on the fluctuating
        complexes without mirror states or broken gauge invariance; or
  \item the continuum high-curvature channel loses norm or retains an
        inaccessible final factor.
\end{enumerate}

\section{Discussion}

CCQG proposes a specific hierarchy of concepts:
\begin{equation}
\boxed{
\begin{gathered}
\text{quantum states and operator algebras}\\
+\ \text{causal composition}\\
+\ \text{finite cut capacity}\\[2mm]
\Downarrow\quad\text{nontrivial refinement fixed point}\quad\Downarrow\\[2mm]
\text{Lorentzian geometry, generalized entropy, and GR}.
\end{gathered}}
\end{equation}
The bridge between quantum theory and gravity is therefore not the scalar
entropy of a state.  It is the algebraic capacity of causal cuts together
with modular response and unitary composition.  Entropy is the
thermodynamic shadow of that structure.

The construction makes three conceptual commitments.  First, spacetime
is not quantized as a pre-existing continuum; a metric is a coherent
phase of nongeometric relational degrees of freedom.  Second, the area
term is a central boundary charge, while statistical entropy depends on
the state occupying the corresponding sector.  Third, singularity
resolution means continuation of relational quantum amplitudes, not
necessarily continuation of a classical manifold.

The conditional infrared argument is comparatively rigid: once an
anomaly-free, local, Lorentz-invariant phase has a single massless
spin-two field and the modular area relation, Einstein gravity is the
expected two-derivative closure.  The genuinely difficult problem is not
guessing the infrared equation.  It is proving that a finite
background-free quantum system lies in that universality class.

\appendix

\section{Notation and scales}

\begin{center}
\begin{tabular}{ll}
\toprule
Symbol & Meaning\\
\midrule
\(\Lambda=(N,Q)\) & finite graph and capacity regulator\\
\(\Gamma,\gamma\) & relational cut or codimension-two screen\\
\(q_f\) & finite local Hilbert-space capacity label\\
\(s_f=\log q_f\) & additive capacity charge\\
\(a_\star\) & microscopic area per nat\\
\(\ell_\star\) & microscopic crossover length\\
\(\lambda_{\mathrm{cc}}\) & cosmological constant\\
\(\widehat H_\Lambda\) & regulated local-move Hamiltonian\\
\(\eta\) & continuum rigging map\\
\(P_{\mathrm{code}}\) & semiclassical code-subspace projector\\
\bottomrule
\end{tabular}
\end{center}

We use \(c=k_{\mathrm B}=1\) but retain \(\hbar\).  The convention
\(a_\star=4G\hbar\) makes capacity measured in nats.  If capacity is
measured in bits, \(a_\star\) acquires a factor of \(\ln 2\).

\section{Minimal numerical program}

Even before an analytic fixed-point proof, the framework admits
discriminating numerical tests:
\begin{enumerate}
  \item enumerate small causal complexes and construct sparse
        \(\widehat H_\Lambda\);
  \item verify exact finite-cutoff unitarity and automorphism invariance;
  \item coarse-grain two-point functions of geometric observables;
  \item measure spectral dimension and the number of gapless modes;
  \item test discrete constraint commutators on coherent states;
  \item measure relative entropy across nested cuts and compare it with
        the capacity-area operator; and
  \item evolve collapse-like states through the no-geometry sector and
        track mutual information between the remaining compact region
        and outgoing degrees of freedom.
\end{enumerate}
Failure at steps 3--5 would reject a candidate move Hamiltonian long
before a full Standard Model or black-hole calculation is attempted.

\begin{thebibliography}{99}

\bibitem{Bekenstein1973}
J.~D. Bekenstein,
``Black holes and entropy,''
\emph{Phys. Rev. D} \textbf{7}, 2333 (1973).
\doi{10.1103/PhysRevD.7.2333}

\bibitem{Hawking1975}
S.~W. Hawking,
``Particle creation by black holes,''
\emph{Commun. Math. Phys.} \textbf{43}, 199 (1975);
Erratum \textbf{46}, 206 (1976).
\doi{10.1007/BF02345020}

\bibitem{Regge1961}
T. Regge,
``General relativity without coordinates,''
\emph{Nuovo Cim.} \textbf{19}, 558 (1961).
\doi{10.1007/BF02733251}

\bibitem{Markopoulou2000}
F. Markopoulou,
``Quantum causal histories,''
\emph{Class. Quantum Grav.} \textbf{17}, 2059 (2000),
\href{https://arxiv.org/abs/hep-th/9904009}{arXiv:hep-th/9904009}.

\bibitem{Hawkins2003}
E. Hawkins, F. Markopoulou, and H. Sahlmann,
``Evolution in quantum causal histories,''
\emph{Class. Quantum Grav.} \textbf{20}, 3839 (2003),
\href{https://arxiv.org/abs/hep-th/0302111}{arXiv:hep-th/0302111}.

\bibitem{RovelliVidotto2015}
C. Rovelli and F. Vidotto,
\emph{Covariant Loop Quantum Gravity},
Cambridge University Press (2015).

\bibitem{Bruno2026}
M. Bruno, E. Colafranceschi, F.~M. Mele, and C. Rovelli,
``The structure of the continuum limit of spin foams,''
\href{https://arxiv.org/abs/2603.16999}{arXiv:2603.16999} (2026).

\bibitem{Jacobson1995}
T. Jacobson,
``Thermodynamics of spacetime: The Einstein equation of state,''
\emph{Phys. Rev. Lett.} \textbf{75}, 1260 (1995),
\href{https://arxiv.org/abs/gr-qc/9504004}{arXiv:gr-qc/9504004}.

\bibitem{Jacobson2015}
T. Jacobson,
``Entanglement equilibrium and the Einstein equation,''
\emph{Phys. Rev. Lett.} \textbf{116}, 201101 (2016),
\href{https://arxiv.org/abs/1505.04753}{arXiv:1505.04753}.

\bibitem{HKT1976}
S.~A. Hojman, K. Kucha\v{r}, and C. Teitelboim,
``Geometrodynamics regained,''
\emph{Ann. Phys.} \textbf{96}, 88 (1976).
\doi{10.1016/0003-4916(76)90112-3}

\bibitem{Weinberg1964}
S. Weinberg,
``Photons and gravitons in \(S\)-matrix theory: Derivation of charge
conservation and equality of gravitational and inertial mass,''
\emph{Phys. Rev.} \textbf{135}, B1049 (1964).
\doi{10.1103/PhysRev.135.B1049}

\bibitem{Deser1970}
S. Deser,
``Self-interaction and gauge invariance,''
\emph{Gen. Rel. Grav.} \textbf{1}, 9 (1970),
\href{https://arxiv.org/abs/gr-qc/0411023}{arXiv:gr-qc/0411023}.

\bibitem{WeinbergWitten1980}
S. Weinberg and E. Witten,
``Limits on massless particles,''
\emph{Phys. Lett. B} \textbf{96}, 59 (1980).
\doi{10.1016/0370-2693(80)90212-9}

\bibitem{Wald1993}
R.~M. Wald,
``Black hole entropy is the Noether charge,''
\emph{Phys. Rev. D} \textbf{48}, R3427 (1993),
\href{https://arxiv.org/abs/gr-qc/9307038}{arXiv:gr-qc/9307038}.

\bibitem{Bueno2017}
P. Bueno, V.~S. Min, A.~J. Speranza, and M.~R. Visser,
``Entanglement equilibrium for higher order gravity,''
\emph{Phys. Rev. D} \textbf{95}, 046003 (2017),
\href{https://arxiv.org/abs/1612.04374}{arXiv:1612.04374}.

\bibitem{ADH2015}
A. Almheiri, X. Dong, and D. Harlow,
``Bulk locality and quantum error correction in AdS/CFT,''
\emph{JHEP} \textbf{04}, 163 (2015),
\href{https://arxiv.org/abs/1411.7041}{arXiv:1411.7041}.

\bibitem{JLMS2016}
D. Jafferis, A. Lewkowycz, J. Maldacena, and S.~J. Suh,
``Relative entropy equals bulk relative entropy,''
\emph{JHEP} \textbf{06}, 004 (2016),
\href{https://arxiv.org/abs/1512.06431}{arXiv:1512.06431}.

\bibitem{Pastawski2015}
F. Pastawski, B. Yoshida, D. Harlow, and J. Preskill,
``Holographic quantum error-correcting codes: Toy models for the
bulk/boundary correspondence,''
\emph{JHEP} \textbf{06}, 149 (2015),
\href{https://arxiv.org/abs/1503.06237}{arXiv:1503.06237}.

\bibitem{EngelhardtWall2015}
N. Engelhardt and A.~C. Wall,
``Quantum extremal surfaces: Holographic entanglement entropy beyond the
classical regime,''
\emph{JHEP} \textbf{01}, 073 (2015),
\href{https://arxiv.org/abs/1408.3203}{arXiv:1408.3203}.

\bibitem{Penington2020}
G. Penington,
``Entanglement wedge reconstruction and the information paradox,''
\emph{JHEP} \textbf{09}, 002 (2020),
\href{https://arxiv.org/abs/1905.08255}{arXiv:1905.08255}.

\bibitem{Almheiri2020}
A. Almheiri, N. Engelhardt, D. Marolf, and H. Maxfield,
``The entropy of bulk quantum fields and the entanglement wedge of an
evaporating black hole,''
\emph{JHEP} \textbf{12}, 063 (2019),
\href{https://arxiv.org/abs/1905.08762}{arXiv:1905.08762}.

\bibitem{GinspargWilson1982}
P.~H. Ginsparg and K.~G. Wilson,
``A remnant of chiral symmetry on the lattice,''
\emph{Phys. Rev. D} \textbf{25}, 2649 (1982).
\doi{10.1103/PhysRevD.25.2649}

\bibitem{Luscher1999}
M. L\"uscher,
``Abelian chiral gauge theories on the lattice with exact gauge
invariance,''
\emph{Nucl. Phys. B} \textbf{549}, 295 (1999),
\href{https://arxiv.org/abs/hep-lat/9811032}{arXiv:hep-lat/9811032}.

\end{thebibliography}

\end{document}
